\documentclass[11pt,a4paper]{article}
\usepackage[utf8]{inputenc}
\usepackage[T1]{fontenc}
\usepackage[spanish,es-nodecimaldot]{babel}
\usepackage{amsmath,amssymb,mathtools}
\usepackage{booktabs,longtable,array}
\usepackage{geometry}
\usepackage{enumitem}
\usepackage{hyperref}
\usepackage{xcolor}
\usepackage{listings}
\usepackage{microtype}
\geometry{margin=2.35cm}
\definecolor{codegray}{gray}{0.96}
\lstset{basicstyle=\ttfamily\small,frame=single,backgroundcolor=\color{codegray},breaklines=true,columns=fullflexible,showstringspaces=false}
\newcommand{\Sargon}{\textsc{Sargon}}
\newcommand{\ValueNet}{\textsc{ValueNet}}
\newcommand{\PieceNet}{\textsc{PieceNet}}
\newcommand{\DepthNet}{\textsc{DepthNet}}
\newcommand{\Model}{\mathcal{M}}
\title{\textbf{Especificación del Motor Sargon Adaptativo}\\Aprendizaje neuronal de valor posicional, valor contextual de piezas y profundidad de análisis}
\author{Especificación de trabajo}
\date{\today}
\begin{document}
\maketitle
\begin{abstract}
Este documento especifica una evolución gradual del algoritmo \Sargon{} mediante tres redes neuronales especializadas: una red de valor posicional, una red de valor contextual de las piezas y una red de profundidad de análisis. La arquitectura se diseña con dos prioridades: preservar el comportamiento de \Sargon{} como línea base cuando no existe aprendizaje efectivo, y hacer que todo progreso sea observable, medible y auditable.

La especificación define la interacción entre las fuentes de juego UCI, \texttt{learn}, workers y archivos PGN; establece un mecanismo uniforme de backpressure por productor; fija la coherencia de un único modelo durante toda una partida; determina cómo se aprende a partir de análisis realizados, resultados de partida y jugadas expertas de PGN; y establece las métricas mínimas que deben registrarse para demostrar si el sistema aprende realmente.

El objetivo explícito es evitar arquitecturas opacas en las que el volumen de datos, la concurrencia o la complejidad de la infraestructura oculten un aprendizaje débil o inexistente.
\end{abstract}
\tableofcontents
\newpage

\section{Propósito y alcance}
El propósito de esta especificación es definir un motor de ajedrez adaptativo basado en \Sargon{} que:
\begin{enumerate}
\item utilice \Sargon{} clásico como conocimiento inicial;
\item introduzca aprendizaje exclusivamente mediante redes neuronales;
\item mantenga separadas tres capacidades aprendidas: valor posicional, valor contextual de las piezas y profundidad necesaria de análisis;
\item pueda aprender de partidas UCI, partidas generadas mediante \texttt{learn}, partidas producidas por workers y archivos PGN;
\item trate todo PGN importado como experiencia experta;
\item mantenga un modelo coherente durante una partida completa;
\item impida que cualquier productor comience una nueva partida hasta que la anterior haya sido incorporada al aprendizaje;
\item disponga de métricas suficientes para determinar si cada red aprende, cuánto aprende y si ese aprendizaje mejora el comportamiento del motor.
\end{enumerate}

Queda fuera del alcance de esta versión:
\begin{itemize}
\item aprendizaje directo mediante tablas persistentes por posición exacta;
\item uso del estado histórico de una posición como conocimiento operativo principal;
\item redes de política separadas;
\item MCTS o variantes tipo AlphaZero;
\item actualización directa de redes desde UCI, workers o \texttt{learn};
\item escritura concurrente sobre los parámetros del modelo.
\end{itemize}

\section{Principios de diseño}
\subsection{Sargon como línea base}
La primera propiedad requerida es:
\[
\boxed{\text{sin aprendizaje efectivo}\Rightarrow\text{\Sargon{} clásico}}
\]
Sea $V_S(s)$ la evaluación clásica de \Sargon{} para una posición $s$. El sistema debe comenzar cumpliendo
\[
V_{\mathrm{eff}}(s)=V_S(s).
\]
Análogamente, el valor contextual inicial de cada pieza debe coincidir con el valor clásico y la profundidad inicial con la profundidad base de \Sargon{}.

\subsection{Aprendizaje neuronal exclusivo}
Toda adaptación se realizará mediante tres redes neuronales:
\[
N_V,\qquad N_P,\qquad N_D.
\]
Sus tamaños serán configurables para permitir experimentación controlada.

\subsection{Separación de responsabilidades}
Se establece la regla:
\[
\boxed{\text{los productores generan experiencia; sólo el learner modifica el modelo}}
\]
Los productores son UCI, \texttt{learn}, cada worker y el importador PGN. Ninguno actualiza pesos neuronales directamente.

\subsection{Coherencia por partida}
Toda partida debe jugarse utilizando un único snapshot del modelo:
\[
\Model_k.
\]
Una partida que comienza con $\Model_k$ termina con $\Model_k$. No se permite reemplazar el modelo en mitad de la partida.

\subsection{Backpressure por productor}
Cada productor queda bloqueado para iniciar una nueva partida hasta recibir confirmación de que su partida anterior fue incorporada al aprendizaje. La espera es independiente por productor:
\[
\boxed{G_i\rightarrow\text{learn}(G_i)\rightarrow ACK(G_i)\rightarrow G_{i+1}}
\]

\section{Modelo global}
Una versión del conocimiento operativo se define como
\[
\Model_k=\{N_V^{(k)},N_P^{(k)},N_D^{(k)},\Theta_k,\mathcal{H}_k\},
\]
donde $\Theta_k$ contiene parámetros de moderación y $\mathcal{H}_k$ metadata, métricas y versión. Cada partida almacena la versión de modelo que la produjo.

\section{Red de valor posicional}
\subsection{Objetivo}
La \ValueNet{} aprende la corrección posicional que \Sargon{} no representa adecuadamente:
\[
N_V(x_V(s))=\Delta V_{\mathrm{pos}}(s).
\]
La evaluación efectiva será
\[
V_{\mathrm{eff}}(s)=V_S(s)+\mu_V\Delta V_{\mathrm{pos}}(s)+\Delta V_{\mathrm{piece}}(s),
\]
donde $\mu_V$ es un moderador.

\subsection{Entrada}
$x_V(s)$ deberá contener una representación normalizada y versionada de, al menos: diferencia material, movilidad, control del centro, desarrollo, seguridad del rey, estructura de peones, peones pasados, actividad de piezas, espacio, iniciativa, lado al turno y fase de partida.

\subsection{Arquitectura configurable}
Ejemplos:
\begin{lstlisting}
--value-net 32
--value-net 64,32
--value-net 128,64,32
--value-net 256,128,64
\end{lstlisting}
La arquitectura usada se persistirá con el modelo.

\subsection{Target de entrenamiento}
Para una posición $s_t$, sean $V_S(s_t)$ la evaluación estática clásica y $V_{d_t}(s_t)$ el resultado de la búsqueda realmente realizada a profundidad $d_t$. La búsqueda proporciona
\[
Y_t^{\mathrm{search}}=V_{d_t}(s_t).
\]
El resultado posterior produce un retorno temporal $Y_t^{\mathrm{TD}}$. El target combinado es
\[
Y_t^V=\alpha_tY_t^{\mathrm{search}}+(1-\alpha_t)Y_t^{\mathrm{TD}}.
\]
La red aprende la corrección residual
\[
\boxed{T_t^V=Y_t^V-V_S(s_t)-\Delta V_{\mathrm{piece}}(s_t)}.
\]
La pérdida básica es
\[
L_V=\frac1B\sum_{t=1}^{B}\left(N_V(x_V(s_t))-T_t^V\right)^2.
\]

\subsection{Retroalimentación por resultado}
No se asignará directamente el resultado final a todas las posiciones. Sea
\[
z\in\{-1,0,+1\}.
\]
Se utilizará aprendizaje temporal con
\[
\delta_t=r_{t+1}+\gamma V_{\mathrm{eff}}(s_{t+1})-V_{\mathrm{eff}}(s_t),
\]
y TD$(\lambda)$ para controlar la propagación de consecuencias tardías hacia posiciones anteriores. El retorno se calcula una vez terminada la partida.

\section{Red de valor contextual de piezas}
\subsection{Objetivo}
La \PieceNet{} aprende cuánto difiere el valor efectivo de una pieza de su valor clásico en función de la posición:
\[
N_P(x_P(s))=(\Delta P,\Delta N,\Delta B,\Delta R,\Delta Q).
\]
Los valores efectivos serán
\[
v_X(s)=v_X^{(0)}+\mu_P\Delta X(s),\qquad X\in\{P,N,B,R,Q\}.
\]

\subsection{Separación semántica}
La \PieceNet{} sólo podrá intervenir en términos asociados directamente con valor contextual de piezas. La evaluación se descompone como
\[
V_{\mathrm{eff}}(s)=V_S(s)+\Delta V_{\mathrm{piece}}(s)+\Delta V_{\mathrm{pos}}(s).
\]
Así se evita que \ValueNet{} y \PieceNet{} aprendan de manera indistinguible el mismo fenómeno.

\subsection{Entrada}
$x_P(s)$ contendrá como mínimo: tipo y cantidad de piezas, movilidad por tipo, centralización, control de casillas, puestos avanzados, columnas abiertas y semiabiertas, diagonales activas, actividad relativa, estructura de peones, fase de partida, seguridad de reyes y relaciones pieza-estructura.

\subsection{Arquitectura configurable}
\begin{lstlisting}
--piece-net 32
--piece-net 64,32
--piece-net 128,64
\end{lstlisting}
La salida tendrá siempre cinco valores.

\subsection{Aprendizaje}
La \PieceNet{} se entrena a partir del mismo target global de valor, pero el gradiente queda restringido a la contribución material contextual. El error global es
\[
E_t=Y_t^V-V_{\mathrm{eff}}(s_t).
\]

\section{Red de profundidad}
\subsection{Objetivo}
La \DepthNet{} aprende cuánto análisis adicional necesita una posición futura:
\[
N_D(x_D(s))=q(s).
\]
La profundidad aplicada es
\[
d_{\mathrm{eff}}(s)=\operatorname{clip}(d_0+\Delta d(s),d_{\min},d_{\max}).
\]

\subsection{Moderadores obligatorios}
La salida neuronal se aplica siempre moderada:
\[
\Delta d(s)=\operatorname{round}(\mu_DN_D(x_D(s))).
\]
Inicialmente
\[
0\le\mu_D\le1,
\qquad
0\le\Delta d(s)\le\Delta d_{\max}.
\]
En la etapa inicial la red puede agregar profundidad, pero no reducir la profundidad base.

\subsection{Target de profundidad}
Durante una búsqueda se conservan
\[
V_1,V_2,\ldots,V_D.
\]
Sea $V_D$ el valor a la mayor profundidad alcanzada. Se define
\[
d^*=\min\{d:|V_d-V_D|<\epsilon\ \text{y la estabilidad se mantiene}\}.
\]
La estabilidad deberá mantenerse durante un número configurable de profundidades posteriores. El target es
\[
T^D=d^*-d_0,
\]
y la pérdida básica
\[
L_D=\frac1B\sum\left(N_D(x_D(s))-T^D\right)^2.
\]

\subsection{Aplicación futura}
El aprendizaje de profundidad no se guarda como regla por posición exacta. Queda codificado en $N_D$. Cuando una posición futura presenta características similares, la red puede solicitar análisis adicional antes de que \Sargon{} descubra tardíamente su dificultad.

\section{Fuentes de experiencia}
\subsection{UCI}
UCI utiliza un snapshot fijo durante toda la partida. Al finalizar genera experiencia, la envía al learner, queda bloqueado y sólo acepta una nueva partida después del ACK.

\subsection{Modo learn}
\texttt{learn} sigue exactamente el mismo protocolo:
\[
\boxed{\texttt{learn game}\rightarrow\text{learner}\rightarrow ACK\rightarrow\text{next learn game}}.
\]

\subsection{Workers}
Cada worker es un productor independiente. El worker $i$ ejecuta
\[
G_i^k\rightarrow\text{learner}\rightarrow ACK_i\rightarrow G_i^{k+1}.
\]
El ACK de un worker no bloquea a otros. Cada worker utiliza el último modelo disponible al iniciar su partida, mantiene ese modelo hasta el final, no entrena redes y produce experiencia adversarial.

\subsection{PGN}
Todo PGN importado se considera fuente experta:
\[
\boxed{\text{PGN importado}=\text{experiencia experta}}.
\]
No existirán categorías separadas de PGN genérico, maestro o engine.

\section{Aprendizaje desde PGN}
\subsection{Preferencia experta}
Para una posición $s_t$, sea $a_t^*$ la jugada del PGN y $a_1,\ldots,a_n$ las jugadas legales. El sistema incorpora
\[
Q(s_t,a_t^*)\gtrsim Q(s_t,a_j),\qquad a_j\ne a_t^*.
\]
No se exige que $a_t^*$ sea la única jugada buena; se trata como evidencia experta de preferencia.

\subsection{Procesamiento}
Al importar una partida se reconstruye cada posición, se extraen features para las tres redes, se ejecuta \Sargon{} con el modelo vigente, se analiza la jugada experta, se comparan alternativas relevantes, se incorpora el resultado final y se generan targets para valor, piezas y profundidad.

\subsection{Pérdida de preferencia}
Se utilizará una pérdida de ranking
\[
L_{\mathrm{PGN}}=\sum_{a\ne a^*}\max\left(0,m-[Q(s,a^*)-Q(s,a)]\right),
\]
con margen $m$ configurable. La experiencia PGN tendrá un peso superior al de experiencia autocreada:
\[
\omega_{\mathrm{PGN}}>1.
\]

\section{Representación de una partida aprendible}
Cada partida producirá un objeto lógico:
\begin{lstlisting}
GameExperience
    game_id
    source
    model_version
    result
    positions[]
\end{lstlisting}
Cada posición contendrá:
\begin{lstlisting}
PositionExperience
    ply
    side_to_move
    features_value
    features_piece
    features_depth
    sargon_static_value
    effective_value
    search_depth
    search_values_by_depth[]
    best_move
    candidate_moves[]
        move
        value
        depth
    played_move
    expert_move      # si source=PGN
\end{lstlisting}
Al terminar la partida se calculan los targets TD, de valor y de profundidad.

\section{Ciclo de aprendizaje de una partida}
\begin{verbatim}
productor toma Model Mk
        |
        v
 juega partida completa
        |
        +-- posicion 1: evaluacion + busqueda
        +-- posicion 2: evaluacion + busqueda
        +-- ...
        +-- posicion N
        |
        v
   resultado final
        |
        v
 GameExperience completo
        |
        v
      learner
        |
        +-- calcula retornos TD(lambda)
        +-- genera targets ValueNet
        +-- genera targets PieceNet
        +-- genera targets DepthNet
        +-- aplica preferencia PGN cuando corresponda
        +-- actualiza redes
        |
        v
    publica Model Mk+1
        |
        v
     ACK(game_id)
\end{verbatim}

\section{Definición de partida aprendida}
Una partida se considera aprendida cuando:
\begin{enumerate}
\item su resultado final fue registrado;
\item sus posiciones fueron convertidas en ejemplos;
\item se calcularon los targets temporales;
\item se aplicaron los targets de análisis;
\item se ejecutó al menos una actualización de todas las redes habilitadas;
\item el modelo resultante fue publicado;
\item la experiencia quedó persistida;
\item se emitió el ACK al productor.
\end{enumerate}
No se requiere entrenar múltiples epochs antes del ACK. La experiencia queda disponible para replay posterior.

\section{Backpressure}
La regla uniforme para UCI, \texttt{learn} y cada worker es
\[
\boxed{\text{máximo una partida pendiente por productor}}.
\]
Con $W$ workers, un UCI y un \texttt{learn}, el máximo teórico de partidas finalizadas pendientes es aproximadamente $W+2$. No existe una barrera global: el bloqueo es local a cada productor.

\section{Snapshots y publicación de modelos}
Al comenzar una partida:
\[
local\_model\leftarrow current\_model.
\]
Toda la partida utiliza ese objeto. El learner no modifica el modelo actual en el lugar: construye $\Model_{k+1}$ y lo publica atómicamente. Los productores que ya juegan conservan $\Model_k$; los que comienzan una nueva partida después del ACK toman la versión más reciente.

\section{Persistencia}
\subsection{Objetivos}
La persistencia sirve para guardar modelos, experiencia, replay, auditoría, rollback y reconstrucción de métricas. No es el mecanismo primario de decisión durante el juego.

\subsection{Modelo}
Debe persistirse, como mínimo:
\begin{lstlisting}
model_version
created_at
value_net_arch
piece_net_arch
depth_net_arch
value_net_weights
piece_net_weights
depth_net_weights
moderators
training_counters
metrics
\end{lstlisting}

\subsection{Experiencia}
\begin{lstlisting}
game_id
source
model_version
result
positions
targets
learned_at
\end{lstlisting}

\subsection{Reuso del conocimiento}
No se utilizará una posición histórica exacta como fuente primaria de decisión futura. El conocimiento reutilizado se encuentra en $N_V,N_P,N_D$:
\[
\boxed{\text{experiencia histórica}\rightarrow\text{pesos neuronales}\rightarrow\text{generalización futura}}.
\]

\section{Moderadores}
Cada red tendrá un moderador independiente:
\[
\mu_V,\qquad\mu_P,\qquad\mu_D.
\]
La configuración inicial será
\[
\mu_V=\mu_P=\mu_D=0.
\]
Ejemplo de interfaz:
\begin{lstlisting}
--value-influence 0.25
--piece-influence 0.25
--depth-influence 0.20
\end{lstlisting}
Esto permite distinguir capacidad de aprendizaje de influencia efectiva sobre el juego.

\section{Métricas obligatorias}
La auditabilidad es un requisito funcional.

\subsection{Métricas por red}
Para cada red se registrará: ejemplos vistos, loss de entrenamiento, loss de validación, media y desviación de salida, magnitud media del gradiente, cambio medio de parámetros, influencia efectiva y composición por fuente.

\subsection{ValueNet}
Se comparará
\[
MAE(V_S,Y^V)
\]
con
\[
MAE(V_{\mathrm{eff}},Y^V).
\]
La mejora es
\[
\Delta MAE=MAE(V_S,Y^V)-MAE(V_{\mathrm{eff}},Y^V).
\]
Reducir sólo el loss interno no es evidencia suficiente.

\subsection{PieceNet}
Se medirán corrección media por pieza, distribución por fase, contribución al error total, estabilidad temporal y magnitud relativa respecto del valor base.

\subsection{DepthNet}
Se medirá
\[
MAE(d_{\mathrm{pred}},d^*)
\]
y además profundidad media solicitada y efectiva, porcentaje de posiciones extendidas, nodos adicionales, error de valor evitado y costo computacional por mejora.

\subsection{Coincidencia con PGN}
Sobre validación PGN no usada para entrenamiento:
\[
A_{\mathrm{top1}}=P(a_{\mathrm{engine}}=a_{\mathrm{expert}}),
\]
\[
A_{\mathrm{topK}}=P(a_{\mathrm{expert}}\in\mathrm{TopK}).
\]

\subsection{Impacto sobre decisiones}
Se registrarán la jugada base $a_S$ y la jugada con redes $a_L$, midiendo
\[
P(a_S\ne a_L).
\]

\subsection{Resultados por generación}
Cada versión registrará cantidad de partidas, victorias, derrotas, tablas, fuente, modelo adversario cuando corresponda, duración media, nodos por jugada, profundidad media y tiempo por jugada.

\section{Validación y conjuntos de prueba}
Deben existir conjuntos congelados de validación: posiciones PGN no usadas en entrenamiento, posiciones tácticas y posicionales de referencia, partidas reproducibles con seeds fijas y enfrentamientos entre generaciones de modelos. No se permitirá evaluar progreso únicamente sobre datos de entrenamiento.

\section{Criterios mínimos para afirmar que una red aprende}
\subsection{ValueNet}
Debe disminuir el error de validación, mejorar $MAE(V_{\mathrm{eff}},Y^V)$ frente a \Sargon{}, generalizar a datos no vistos y no degradar sistemáticamente el resultado del motor.

\subsection{PieceNet}
Sus correcciones deben ser estables, mejorar el target global, generalizar y no colapsar a correcciones extremas.

\subsection{DepthNet}
Debe predecir $d^*$ mejor que una profundidad constante, reducir error de evaluación mediante extensiones, justificar el costo de nodos adicionales y permanecer dentro de límites de profundidad.

\section{Auditoría obligatoria}
Cada entrenamiento deberá permitir reconstruir: partida origen, versión del modelo, evaluación Sargon, correcciones neuronales, profundidad usada, target posterior al resultado, pérdidas producidas y modelo resultante. Cada modelo debe ser trazable a
\[
\text{modelo anterior}+\text{partidas aprendidas}+\text{configuración}+\text{seed}
\]
cuando sea posible.

\section{Logging de consola}
El sistema deberá reportar de forma compacta:
\begin{lstlisting}
[game-start] source=worker-3 model=42
[game-end] source=worker-3 game=918 result=0-1 plies=83
[learn-start] game=918 model_in=42
[learn-value] samples=83 loss=...
[learn-piece] samples=83 loss=...
[learn-depth] samples=... loss=...
[learn-end] game=918 model_out=43 elapsed=...
[ack] source=worker-3 game=918
\end{lstlisting}
Cada productor reportará cuándo espera ACK y cuándo queda habilitado para una nueva partida.

\section{Configuración mínima}
La CLI deberá permitir, como mínimo:
\begin{lstlisting}
--value-net <layers>
--piece-net <layers>
--depth-net <layers>
--value-influence <float>
--piece-influence <float>
--depth-influence <float>
--depth-base <int>
--depth-min <int>
--depth-max <int>
--depth-max-extra <int>
--gamma <float>
--lambda <float>
--pgn-weight <float>
--learning-rate-value <float>
--learning-rate-piece <float>
--learning-rate-depth <float>
\end{lstlisting}
La configuración exacta se almacenará en cada modelo.

\section{Secuencia de activación}
Aunque la infraestructura soporte las tres redes desde el comienzo, la activación funcional será escalonada.

\subsection{Etapa 0}
\[
\mu_V=\mu_P=\mu_D=0.
\]
Debe coincidir con \Sargon{} base.

\subsection{Etapa 1}
\[
\mu_V>0,\qquad\mu_P=0,\qquad\mu_D=0.
\]
Se entrena y audita \ValueNet{}.

\subsection{Etapa 2}
\[
\mu_V>0,\qquad\mu_P>0,\qquad\mu_D=0.
\]
Se incorpora \PieceNet{}.

\subsection{Etapa 3}
Se incorpora \DepthNet{} con moderación y sólo extensiones:
\[
\mu_D>0,\qquad\Delta d\ge0.
\]

\section{Arquitectura general}
\begin{verbatim}
                       current Model Mk
                  NV + NP + ND + Sargon
                            |
      +---------------------+----------------------+
      |                     |                      |
     UCI                worker 1...N             learn
      |                     |                      |
      +----------- GameExperience ----------------+
                            |
                            v
                         Learner
                            |
        +-------------------+-------------------+
        |                   |                   |
      ValueNet            PieceNet            DepthNet
        |                   |                   |
        +-------------------+-------------------+
                            |
                            v
                        Model Mk+1
                            |
              +-------------+-------------+
              |             |             |
             ACK           ACK           ACK
             UCI         worker i       learn
\end{verbatim}
PGN entra como experiencia experta:
\begin{verbatim}
PGN ---> parser ---> GameExperience(expert) ---> Learner
\end{verbatim}

\section{Ciclo conceptual del conocimiento}
Para valor:
\[
\boxed{\text{posición}\rightarrow\text{evaluación Sargon}\rightarrow\text{análisis}\rightarrow\text{target provisional}\rightarrow\text{resultado}\rightarrow\text{target temporal}\rightarrow\text{redes}\rightarrow\text{modelo futuro}}
\]
Para PGN:
\[
\boxed{\text{posición}\rightarrow\text{jugada experta}\rightarrow\text{comparación con alternativas}\rightarrow\text{pérdida de preferencia}\rightarrow\text{redes}}
\]
Para profundidad:
\[
\boxed{V_1,\ldots,V_D\rightarrow d^*\rightarrow N_D\rightarrow\text{profundidad futura moderada}}
\]

\section{Invariantes arquitectónicos}
Se establecen como invariantes:
\begin{enumerate}
\item sólo existe un escritor lógico del modelo;
\item ningún productor modifica pesos;
\item una partida usa un único modelo;
\item una fuente no inicia una nueva partida antes del ACK;
\item cada worker posee su propia espera independiente;
\item todos los modelos son versionados;
\item todas las experiencias indican el modelo que las produjo;
\item todos los PGN son experiencia experta;
\item el conocimiento operativo reside en las redes, no en tablas por posición exacta;
\item la profundidad neuronal siempre está moderada y limitada;
\item ninguna mejora se considera demostrada sin validación fuera de muestra.
\end{enumerate}

\section{Objetivo experimental}
La arquitectura debe permitir responder cuantitativamente:
\begin{enumerate}
\item ¿disminuye el error de valor?;
\item ¿cambia de manera estable la valoración contextual de piezas?;
\item ¿DepthNet identifica posiciones que realmente requieren más análisis?;
\item ¿las redes generalizan a posiciones no vistas?;
\item ¿los PGN aumentan la coincidencia con decisiones expertas?;
\item ¿los workers producen experiencia que mejora generaciones posteriores?;
\item ¿el modelo nuevo supera al anterior en métricas independientes?;
\item ¿el aumento de complejidad produce una mejora proporcional y verificable?
\end{enumerate}
Si estas preguntas no pueden responderse con datos del propio sistema, la implementación no cumple esta especificación.

\section{Conclusión}
La arquitectura queda definida por cuatro ideas fundamentales. Primero, \Sargon{} continúa siendo la línea base y proporciona estructura, evaluación y búsqueda. Segundo, el aprendizaje se concentra en tres redes neuronales especializadas:
\[
\boxed{N_V,\quad N_P,\quad N_D}
\]
para valor posicional, valor contextual de piezas y profundidad. Tercero, toda fuente de juego produce experiencia, pero sólo el learner actualiza el modelo:
\[
\boxed{\text{todos producen; uno aprende}}.
\]
Cuarto, el progreso debe ser demostrable mediante métricas reproducibles y validación independiente.

La especificación prioriza deliberadamente simplicidad, trazabilidad y capacidad de falsar la hipótesis de aprendizaje. El sistema debe poder demostrar no sólo que sus pesos cambian, sino que esos cambios reducen error, mejoran decisiones, utilizan mejor la profundidad y producen generaciones de juego progresivamente superiores.
\end{document}
